\chapter{Classical Game Theory: Extensive Form and Its Interpretation}
\label{sec:p17-rationality-9-1}

An equilibrium calculation and an explanation of human interaction are
different achievements. Albert and Kliemt's classical-game-theory chapter
first specifies a mathematical language, then examines the additional
interpretations needed to use it. Its central object is an extensive game,
not just a payoff matrix. Its principal refinement is subgame perfection;
its final argument questions whether correct predictions of behavior must
come from reasoning about common knowledge
\parencite[chap.~9.1, pp.~529--541]{AlbertKliemt2021ClassicalGameTheory}.
The treatment below preserves that division. Historical references identify
origins and alternatives; they are not empirical evidence that actual players
satisfy the mathematical hypotheses.

\section{Rules, histories and information}

\begin{hbdefinition}{def:p17-rat91-tree}{The finite extensive language}
A finite extensive game is a finite rooted tree. Every node other than
the root has one immediate predecessor, and therefore one history from
the root. A complete history ends at a terminal node, which carries
one real payoff for each player. Nonterminal nodes belong to players;
a player without any such node may still receive a payoff.

Each player's decision nodes are partitioned into information sets
\(H\). The same finite set \(A(H)\) of at least two moves labels the
outgoing branches at every node in \(H\); a label identifies the same
move throughout that set. Changing labels at only one of its nodes
changes the game. Information sets represent what the mover can
distinguish, including possible failure to observe, notice or remember.
The moves specified by the tree exhaust this model's available moves;
an unlisted act is not a rule-breaking strategy within the same game.

The source calls a game \emph{strictly sequential} when every
information set is a singleton; this is the usual perfect-information
case. Its \emph{simultaneous} games have at most one information set
per player and every complete history passes through all those sets.
Games can satisfy neither description. An information set, a history
and a terminal outcome are distinct objects.

Nature may be represented by a mover without a payoff, with specified
probabilities at its nodes. In the numerical examples below Nature
is absent. Unless explicitly discussing absentmindedness, no history
visits an information set twice.
\index{extensive game}\index{information set}\index{perfect information}
\end{hbdefinition}

The primitive mathematical words do not by themselves specify an
application. A player might represent a person, organization, organism,
or one role of a person. Probabilities obey the probability axioms but
need an interpretation; a real payoff need not yet mean money, welfare,
or revealed-preference utility. The source adopts familiar choice
language for readability without making these interpretations automatic.
The formal scope is finite: finitely many terminal histories, each of
finite length \parencite[pp.~529--530]{AlbertKliemt2021ClassicalGameTheory}.

\begin{hbdefinition}{def:p17-rat91-strategies}{Contingent plans and two kinds of randomization}
A pure strategy \(s_i\) assigns one move to every information set of
player \(i\), including sets its own earlier moves would prevent from
being reached. A mixed strategy \(\sigma_i\) is a probability on the
finite set of these pure plans. A behavioral strategy \(b_i\) assigns
to each information set a probability on its available moves.
Behavioral randomization is independent at distinct information sets;
in an explicitly absentminded model, fresh independent draws are made
at repeated visits as well. Mixtures chosen by different players are
independent in this classical presentation.

The source writes a lottery as \(p\circ E\oplus q\circ F\), with
nonnegative weights summing to one. This denotes a probability law,
not an operation on payoffs. A pure plan is a degenerate mixture.
A strategic or normal form records all pure-plan profiles and their
induced terminal payoffs, or Nature-averaged payoffs if Nature is present.
It may give identical rows to different plans and may be shared by
different extensive games. The normal-form action set in
\cref{def:p17-finite-strategic-game} is here the set of entire pure plans,
not the collection of individual tree moves.
\index{strategy!mixed versus behavioral}
\end{hbdefinition}

\begin{hbdefinition}{def:p17-rat91-first-trees}{Source Figures 9.1.1 and 9.1.2 in full}
Player 1 first chooses \(A\) or \(B\). After \(A\), player 2
chooses \(W\) or \(X\); after \(B\), \(Y\) or \(Z\).
After \(AW\), player 1 chooses \(C\) or \(D\).
The five terminal histories and payoffs are
\[
\begin{array}{c|ccccc}
\text{history}&AWC&AWD&AX&BY&BZ\\\hline
\text{payoff}&(5,0)&(4,6)&(2,1)&(0,3)&(3,4).
\end{array}
\]
All information sets are singleton in Figure 9.1.1. Player 1 has
plans \(AC,AD,BC,BD\); player 2 has \(WY,WZ,XY,XZ\).
The profile \((BC,XZ)\) selects \(BZ\), giving \((3,4)\),
although its unplayed instructions still specify \(C\) and \(X\).

In Figure 9.1.2 the two player-2 nodes instead share one information
set with moves \(W,X\). Replace histories \(BY,BZ\) by \(BW,BX\)
with the same payoffs. Player 2 now has only the two plans \(W,X\).
The profile \((AD,X)\) selects \(AX\), giving \((2,1)\).
The complete normal forms, source Table 9.1.1, are
\[
\begin{array}{c|cccc}
 &WY&WZ&XY&XZ\\\hline
 AC&(5,0)&(5,0)&(2,1)&(2,1)\\
 AD&(4,6)&(4,6)&(2,1)&(2,1)\\
 BC&(0,3)&(3,4)&(0,3)&(3,4)\\
 BD&(0,3)&(3,4)&(0,3)&(3,4)
\end{array}
\qquad
\begin{array}{c|cc}
 &W&X\\\hline
 AC&(5,0)&(2,1)\\AD&(4,6)&(2,1)\\
 BC&(0,3)&(3,4)\\BD&(0,3)&(3,4).
\end{array}
\]
\end{hbdefinition}

Repeated rows do not erase the extensive instructions: \(BC\) and
\(BD\) induce the same terminal outcome against every opponent plan,
but prescribe different responses if the last player-1 node is
considered separately. Joining player 2's nodes removes its ability
to condition a move on which first move player 1 made. This is an
information restriction, not a change of player 2's preferences
\parencite[pp.~530--532]{AlbertKliemt2021ClassicalGameTheory}.

\section{Equilibrium and the source's simultaneous games}

\begin{hbdefinition}{def:p17-rat91-equilibrium}{Solution, best reply and strictness}
A solution concept assigns selected profiles to games with a stated
strategy domain. A Nash equilibrium is a profile at which no player
can raise its expected payoff by changing only its own strategy.
The complementary profile remains fixed. This is exactly
\cref{def:p17-nash-equilibrium} on the normal form's pure-plan sets.
It is strict when each player's equilibrium strategy is its only
best reply. A mixed equilibrium cannot be strict if some player
uses more than one pure plan with positive probability.

A pure equilibrium without Nature gives one equilibrium history.
With randomization, the equilibrium path means the set of terminal
histories with positive probability. A zero-probability information
set is off that path; its instruction still belongs to the profile.
\end{hbdefinition}

For a finite game, independent mixed Nash existence is already proved
in \cref{thm:p17-nash-1950,thm:p17-nash-1951}. The support criterion
\cref{thm:p17-best-response-geometry} explains why every positively
weighted pure plan must attain the best-response value. It also shows
why equal payoff among supported plans does not force any particular
physical randomization procedure. Correlated recommendations are
expressly beyond the source's classical independent-mixture domain;
the extension and its different obedience test appear in
\cref{def:p17-correlated-recommendation,thm:p17-product-and-correlation}.

\begin{hbdefinition}{def:p17-rat91-four-games}{All four source Figure 9.1.3 games}
Rows are \(U,D\), columns \(L,R\), and entries list row and column
payoffs. The source's numerical games are
\[
\begin{array}{c|cc}\text{Prisoner's Dilemma}&L&R\\\hline
 U&(3,3)&(0,4)\\D&(4,0)&(1,1)\end{array}
\qquad
\begin{array}{c|cc}\text{Matching Pennies}&L&R\\\hline
 U&(1,-1)&(-1,1)\\D&(-1,1)&(1,-1)\end{array}
\]
\[
\begin{array}{c|cc}\text{Battle of the Sexes}&L&R\\\hline
 U&(4,2)&(0,0)\\D&(0,0)&(2,4)\end{array}
\qquad
\begin{array}{c|cc}\text{unnamed}&L&R\\\hline
 U&(4,4)&(0,3)\\D&(4,0)&(1,1).\end{array}
\]
The Prisoner's Dilemma structure has \(x_3>x_1>x_4>x_2\)
and \(y_2>y_1>y_4>y_3\), where entries are ordered
\(UL,UR,DL,DR\). Matching Pennies has
\(x_1=x_4>x_2=x_3\), with the opposite column ordering.
Battle of the Sexes has \(x_1>x_4>x_2=x_3\) and
\(y_4>y_1>y_2=y_3\). Those ordinal patterns do not by
themselves determine the numerical mixed-equilibrium weights.
\end{hbdefinition}

\begin{hbtheorem}{thm:p17-rat91-calculations}{proved=annex:p17.rat91.calculations}{Complete simultaneous calculations and their local sensitivity}
Let \(p\) be the probability of \(U\), and \(q\) that of \(L\).
For general payoffs \((x_k,y_k)\) in the stated cell order,
\[
 U_1=pqx_1+p(1-q)x_2+(1-p)qx_3+(1-p)(1-q)x_4,
\]
with the analogous expression for \(U_2\) using \(y_k\).
The source examples have these complete equilibrium sets:
\[
\begin{array}{c|c|c}
\text{game}&(p,q)&\text{payoffs}\\\hline
\text{Prisoner's Dilemma}&(0,0)&(1,1)\\
\text{Matching Pennies}&(1/2,1/2)&(0,0)\\
\text{Battle of the Sexes}&(1,1),(0,0),(2/3,1/3)&
 (4,2),(2,4),(4/3,4/3)\\
\text{unnamed}&(0,0)\text{ or }q=1,\ 1/2\leq p\leq1&
 (1,1)\text{ or }(4,4p).
\end{array}
\]
The two pure Battle-of-the-Sexes equilibria and the pure Prisoner's
Dilemma equilibrium are strict. The unnamed \((U,L)\) is not.

Adding 1 only to Matching Pennies' row payoff in \(UL\) changes
the unique equilibrium to \((p,q)=(1/2,2/5)\), with payoffs
\((1/5,0)\). On a maintained nondegenerate completely mixed
two-by-two branch, changes in one player's payoffs change the
opponent's indifference equation, not the equation determining the
first player's mixture. At a boundary, with ties or after a change
of equilibrium branch, that local observation is not a global
invariance theorem.
\end{hbtheorem}

The pure-equilibrium procedure is to mark each player's best responses
in each fixed-opponent column or row and retain cells marked for both.
The Prisoner's Dilemma is exceptional in not requiring a correct forecast
of the opponent's action: \(D\) strictly beats \(U\) and \(R\)
strictly beats \(L\) for every opponent mixture. Yet \((3,3)\)
at \((U,L)\) strictly Pareto improves the equilibrium \((1,1)\).
Rationality for supplied all-inclusive utilities is not collective
efficiency \parencite[pp.~532--533]{AlbertKliemt2021ClassicalGameTheory}.

\section{Subgames, credibility and perfect recall}

\begin{hbdefinition}{def:p17-rat91-subgame}{Subgame and subgame-perfect equilibrium}
A subgame starts at a singleton information set, includes every
descendant and its terminal payoff, and never cuts an information set:
every original information set is either fully included or disjoint.
The full game is its own improper subgame. A behavioral profile is
subgame perfect when its restriction is Nash in every subgame,
including the full game.

Perfect recall means that at any two nodes in one information set,
the player's earlier information sets and its own chosen moves occur
in the same ordered sequence. It rules out both forgetting one's own
earlier move and revisiting the same information set along a history.
It does not mean that opponents' unobserved moves become known.
\index{subgame perfection}\index{perfect recall}
\end{hbdefinition}

\begin{hbtheorem}{thm:p17-rat91-realization}{proved=annex:p17.rat91.realization}{The finite perfect-recall mixed--behavioral bridge}
In a finite perfect-recall game every mixed strategy has a behavioral
strategy with the same probability for every terminal history against
every fixed independent opponent profile. Conversely every behavioral
strategy is represented by a mixture of pure plans with those same
terminal probabilities. Hence the two domains give the same attainable
unilateral-deviation payoffs against a fixed behavioral opponent profile.

For the forward construction let \(r_i(H)\) be the mixed mass of
plans compatible with the player's own remembered history before
\(H\), and \(r_i(H,a)\) that of compatible plans also prescribing
\(a\) there. Where \(r_i(H)>0\), use
\(b_i(a\mid H)=r_i(H,a)/r_i(H)\); where it is zero, specify
any move distribution. These are own-history weights, not unconditional
probabilities of actually reaching \(H\). The annex proves the
telescoping realization identity and its zero-weight case.
\end{hbtheorem}

\begin{hbtheorem}{thm:p17-rat91-spe-existence}{proved=annex:p17.rat91.spe.existence}{The finite existence and recursive solution method on its proper domain}
Every finite perfect-recall game with real terminal payoffs and fixed
Nature probabilities has a behavioral subgame-perfect equilibrium.
Recursively solve proper subgames, replace each maximal proper subgame
by its equilibrium payoff, find a mixed Nash equilibrium of the
remaining finite normal form, behavioralize it, and restore the saved
continuations. Terminal histories are finite, so this recursion terminates.

For a strictly sequential game, start at the last decision nodes and
choose a maximizing move, then work toward the root. This gives a pure
subgame-perfect equilibrium. It is unique if each maximization has a
unique optimal move; with ties, every maximizing selection must remain
available. The perfect-recall condition is essential to the general
behavioral existence statement, not a consequence of finiteness alone.
\end{hbtheorem}

Mixed plans alone do not canonically specify conditional random moves
after an own zero-probability history: the forward ratio is then \(0/0\).
Behavioral profiles retain explicit choices there. This is the reason
for the source's emphasis on behavioral descriptions of off-path play,
not a claim that a mixture contains no assignments to pure off-path
instructions \parencite[pp.~533--535]{AlbertKliemt2021ClassicalGameTheory}.

\begin{hbtheorem}{thm:p17-rat91-tree-solutions}{proved=annex:p17.rat91.tree.solutions}{The two source trees distinguish Nash from perfection}
The singleton-node game of \cref{def:p17-rat91-first-trees} has
pure Nash profiles \((AC,XY),(BC,XZ),(BD,XZ)\), but only
\((BC,XZ)\) is subgame perfect. Backward optimization chooses
\(C,Z,X,B\) in that order and yields payoff \((3,4)\).
The two-player-2-node information-set game has pure Nash profiles
\((BC,X),(BD,X)\), but only \((BC,X)\) is subgame perfect.
Its proper subgame requires \(C\); then \(X\) is strictly dominant
for player 2 and \(B\) is the unique best initial response.
\end{hbtheorem}

\begin{hbdefinition}{def:p17-rat91-outside}{The source's Figure 9.1.4 outside option}
Player 1 chooses exit \(X\), giving \((3,3)\), or start \(S\).
After \(S\), player 2 chooses \(L,R\), followed by player 1's
\(U,D\) without observing that choice. The subgame is the stated
Battle of the Sexes. Player 1's pure plans are \(SU,SD,XU,XD\).
Let \(M=\tfrac23\circ U\oplus\tfrac13\circ D\) and
\(N=\tfrac13\circ L\oplus\tfrac23\circ R\).
The full normal form is
\[
\begin{array}{c|cc}&L&R\\\hline
 SU&(4,2)&(0,0)\\SD&(0,0)&(2,4)\\
 XU&(3,3)&(3,3)\\XD&(3,3)&(3,3).
\end{array}
\]
\end{hbdefinition}

\begin{hbtheorem}{thm:p17-rat91-outside}{proved=annex:p17.rat91.outside}{All subgame-perfect choices and the source's forward contrast}
The three pure Nash profiles are \((SU,L),(XU,R),(XD,R)\).
The complete behavioral subgame-perfect set has exactly three profiles:
\((SU,L),(XD,R),(XM,N)\). Their payoffs are respectively
\((4,2),(3,3),(3,3)\). In the mixed continuation the subgame
payoff is \((4/3,4/3)\), making exit strictly optimal; in \((D,R)\)
the row payoff 2 also makes exit optimal, while in \((U,L)\)
the row payoff 4 makes start optimal.

Under the separate explanatory rule that player 2 excludes player 1's
strictly dominated initial plans \(SD\) and \(SM\), observing
\(S\) excludes continuations \(D\) and \(M\). Among the three
subgame-perfect candidates only \((SU,L)\) survives. This is the
source's worked forward-induction contrast under that extra rule,
not a general definition or existence theorem for forward induction.
\end{hbtheorem}

\begin{hbdefinition}{def:p17-rat91-horse}{The complete source Horse Game, Figure 9.1.5}
Player 1 chooses \(D\) or \(R\). After \(R\), player 2 chooses
\(D\) or \(R\); \(RR\) terminates at \((1,1,1)\).
After either \(D\) by player 1 or \(RD\), player 3 moves at
one shared information set and chooses \(L\) or \(R\):
\[
\begin{array}{c|ccccc}
\text{history}&DL&DR&RDL&RDR&RR\\\hline
\text{payoff}&(0,0,0)&(3,2,2)&(0,0,1)&(4,4,0)&(1,1,1).
\end{array}
\]
No proper subgame exists: starting at player 2's node would cut
player 3's information set.
\end{hbdefinition}

\begin{hbtheorem}{thm:p17-rat91-horse}{proved=annex:p17.rat91.horse}{Perfection does not settle every off-path information set}
Both displayed profiles \((D,R,R)\) and \((R,R,L)\) are Nash
and therefore subgame perfect. In the first, player 2's off-path
\(R\) would give 1, while changing locally to \(D\) would give 4.
In the second, player 3's information set has probability zero.
If \(\mu\) is a proposed probability of its left node, player 3's
conditional payoffs are \(1-\mu\) for \(L\) and \(2\mu\)
for \(R\). Thus \(L\) is optimal when \(\mu\leq1/3\),
and \(R\) when \(\mu\geq1/3\). The equilibrium strategies
do not define \(\mu\) by ordinary conditioning: its denominator
is zero. An additional assessment/refinement must supply such beliefs.
\end{hbtheorem}

\section{From a formal game to rational-choice interpretation}

Albert and Kliemt treat an applied model as both a theory and a
description of the situation to which it is applied. For human
rational-choice applications they interpret payoffs as utilities and
probabilities as physical chances or subjective beliefs. They criticize
an unrestricted common-prior assumption as an added hypothesis, not
something implied by the probability axioms. Their applied discussion
assumes rational players and correct beliefs about opponents' strategies;
the required psychological or institutional explanation of those facts
is not supplied by the game tree itself
\parencite[pp.~535--538]{AlbertKliemt2021ClassicalGameTheory}.

\begin{hbdefinition}{def:p17-rat91-rct}{Preferences, choice dispositions and utility}
A complete transitive weak preference relation on a supplied option
domain has strict preference and indifference as its asymmetric and
symmetric parts. A real function represents it when weak preference
holds exactly when the first utility is at least the second.
The finite-domain representation is \cref{thm:p17-finite-representation}.
In the source's revealed-preference interpretation, strict preference
records a disposition not to choose the inferior option when both
are available. Choices therefore maximize the representing utility;
ties permit any maximizing option.

On this interpretation preferences summarize hypothetical choice,
not the mental reasoning that causes it. They are all-inclusive:
material interests, concern for others, moral motivations and procedural
concerns must already be reflected in the stated option ordering.
Utility is not automatically felt well-being. A higher utility means
preferred under that relation, not greater independently measured welfare.
\end{hbdefinition}

An applied predictor therefore needs additional hypotheses about which
motivations shape the ordering and when players possess the information
and accurate beliefs used in the calculation. The chapter mentions
shared cultural backgrounds, experience and sufficiently important
stakes as possible contextual hypotheses, not universal sufficient
conditions established by an experiment. Different such hypotheses
make different theories using a common mathematical language.

\begin{hbexample}{def:p17-rat91-marriage}{Preferences over a path need not be choices over that path}
Eve chooses whether to accept Adam's proposal. Adam may want the
acceptance history more than the refusal history. It does not follow
that his wanting this can be defined by how he would choose if he
controlled Eve's response: being freely accepted by Eve and selecting
acceptance for her may be different objects of his concern.
If the identity of the chooser matters, an ordering of actual
histories is not automatically a revealed-preference ordering over
histories hypothetically controlled by Adam. If what matters instead
is a separable consequence such as a monetary receipt, preferences
over histories may be inherited from preferences over that consequence.
This is the chapter's consequentialism/separability condition, not a
denial that nonconsequentialist preferences can be modeled with a richer
outcome description \parencite[p.~536]{AlbertKliemt2021ClassicalGameTheory}.
\end{hbexample}

\begin{hbdefinition}{def:p17-rat91-uncertainty}{The source's action--state--consequence table}
An action is controlled by the decision maker, a state describes
uncertain aspects outside that control, and a consequence results
from their combination. In source Table 9.1.2 the complete map is
\[
\begin{array}{c|ccc}&S_1&S_2&S_3\\\hline
 A_1&C_1&C_1&C_1\\A_2&C_1&C_2&C_3\\
 A_3&C_5&C_2&C_1\\A_4&C_4&C_2&C_1.
\end{array}
\]
Each set is mutually exclusive and exhaustive in its own typed role.
Given a state law \((\pi_1,\pi_2,\pi_3)\) and consequence
utilities \(u(C_k)\), expected values are
\[
\begin{split}
 u(C_1),\quad
 \pi_1u(C_1)+\pi_2u(C_2)+\pi_3u(C_3),\\
 \pi_1u(C_5)+\pi_2u(C_2)+\pi_3u(C_1),\quad
 \pi_1u(C_4)+\pi_2u(C_2)+\pi_3u(C_1).
\end{split}
\]
This computation requires both the law and utilities; the table
alone does not rank actions.
\end{hbdefinition}

Under risk, known physical probabilities turn actions into lotteries.
An expected-utility representation requires more than an ordinal
ranking of consequences; the finite lottery axioms and complete proof
are in \cref{thm:p17-finite-vnm}. The source contrasts this with
Savage's subjective approach, in which preferences over a rich act
domain yield both probability and utility. The finite provided-belief
calculation above is not a substitute proof of Savage's representation.
The chapter notes gambling and procedural fairness as potential failures
of the required lottery axioms, not as contradictions within the
representation theorem's stated domain. Its common applied interpretation
allows beliefs to be mental states even while utilities remain
representations of hypothetical choices
\parencite[pp.~536--537]{AlbertKliemt2021ClassicalGameTheory}.

\section{Information, time, memory and prediction}

In the applied reading, a player recognizes its current information
set but need not identify its node. Observing another move requires
some causal transmission and implies that the observed move occurred
earlier. Tree order need not encode physical chronology when moves
are unobserved: the simultaneous matrix can be represented with either
player listed first. The source argues that some Newcomb-type causal
arrangements need more than this standard tree language. That is a
limitation of the proposed representation, not a theorem that every
causal phenomenon is a finite game.

\begin{hbdefinition}{def:p17-rat91-memory}{Both source Figure 9.1.6 memory examples}
In (a), player 1 first chooses \(L,R\). Player 2 observes this and
at the corresponding singleton node chooses \(X,D\). Each \(X\)
ends at \((1,1)\). After \(D\), player 1 chooses \(L,R\)
at a shared information set, having forgotten its first move. The
four continuation payoffs are
\[
 LDL:(3,0),\quad LDR:(0,3),\quad RDL:(0,3),\quad RDR:(3,0).
\]
Player 2's two-letter plans list left-node then right-node moves.
The first and last roles of player 1 can alternatively be modeled
as separate players with the same terminal utility; that changes the
unilateral-deviation object.

In (b), player 1 chooses \(L,R\) at the root; \(R\) ends at
\((3,4)\). After \(L\), player 2 chooses \(W,X\);
\(X\) ends at \((2,1)\). After \(LW\), player 1 moves again
at the \emph{same} information set as at the root, not knowing
whether a previous own move has occurred. Its \(L,R\) now end
at \((5,0),(0,2)\). Thus absentmindedness is stronger than
forgetting a remembered earlier move at a different information set.
\end{hbdefinition}

\begin{hbtheorem}{thm:p17-rat91-memory}{proved=annex:p17.rat91.memory}{Memory calculations and a genuine existence boundary}
In Figure 9.1.6a, \((LL,XX)\), or \((L,XX,L)\) with separate
first/last roles, is an equilibrium with payoff \((1,1)\).
The asserted comparison \((LL,DX)\) is not an equilibrium:
player 2 can increase its payoff from 0 to 1 at the reached left
node by changing \(D\) to \(X\).

In Figure 9.1.6b, fix player 2 at \(W\). A behavioral probability
\(1/2\) of \(L\) at each visit yields terminal probabilities
\(1/2,1/4,1/4\) on root \(R\), \(LWL\), \(LWR\),
and payoffs \((11/4,5/2)\). Mixing equally between the pure
plans always-\(L\) and always-\(R\) instead yields payoffs
\((4,2)\). Equal marginal move weights therefore do not give
realization equivalence in this absentminded domain.

Moreover, there exists a finite absentminded two-player game with
independent mixed Nash equilibria but \emph{no behavioral Nash
equilibrium}. The annex specifies its entire tree and payoffs and
proves nonexistence. Consequently the behavioral existence conclusion
of \cref{thm:p17-rat91-spe-existence} cannot drop perfect recall.
No general theory of Sleeping Beauty credences is inferred from these
examples.
\end{hbtheorem}

\begin{hbdefinition}{def:p17-rat91-blackmail}{Source Figures 9.1.7 and 9.1.8: explicit commitment}
In Blackmail, Eve chooses silence \(S\), giving \((2,2)\), or
extortion \(X\). After \(X\), Adam chooses report \(R\),
giving \((0,0)\), or pay \(P\), giving \((3,1)\).
In the extended game Adam first chooses \(N\), leaving this game
unchanged, or observable commitment \(C\), removing the option
\(P\) after extortion. After \(C\), silence still gives \((2,2)\)
and extortion automatically gives \((0,0)\).
Adam's plan \(CP\) records commitment and paying after the unchosen
\(N\) branch. Eve's plan \(XS\) records extorting after \(N\)
and silence after \(C\).
\end{hbdefinition}

\begin{hbtheorem}{thm:p17-rat91-blackmail}{proved=annex:p17.rat91.blackmail}{Credible prediction is not an implicit commitment move}
The unextended game has pure Nash profiles \((X,P),(S,R)\).
Only \((X,P)\) is subgame perfect: Adam prefers paying to reporting
when extortion occurs, and Eve prefers its correctly anticipated
payoff 3 to silence's 2. In the extended game the unique
subgame-perfect profile is \((CP,XS)\), with payoff \((2,2)\).
Its contingent continuations are \((X,P)\) after \(N\) and
\(S\) after \(C\); Adam then prefers commitment's 2 to 1.
\end{hbtheorem}

The threatened report in \((S,R)\) is only a hypothetical instruction.
Without a commitment move, declaring oneself committed does not change
Adam's later incentives. In noncooperative modeling, a genuine device
for removing options must appear in the game. The theory does not tell
us that such a device exists, whether commitment is physical or mental,
or that willpower is observable. Those are factual premises of the
extended model. The same explicitness condition applies to proposals
to escape the supplied Prisoner's Dilemma: different motivations change
utilities; communication, enforcement or commitment changes the tree.
Neither change is a different solution of the original unaltered game
\parencite[pp.~539--540]{AlbertKliemt2021ClassicalGameTheory}.

Correct beliefs are the source's strong applied premise, including
beliefs about what opponents would do after a history assigned zero
probability. A player's contingent expectations already incorporate
the unexpected earlier move needed to reach that history. In the first
source tree, anticipating \(C\) after surprising \(A\) need not
be deduced from an unqualified belief that nobody ever chooses wrongly.
The source therefore separates correct predictions of continuation
strategies from knowledge of opponents' utilities or rationality.
On that stipulated correct-counterfactual-belief interpretation,
unexpected arrival need not trigger a revision of predictions about
the subsequent moves: those predictions were already conditional
on arriving there. This is a substantive modeling premise, not a
general psychological law of belief revision.
At the Horse Game's unreached nonsingleton information set, however,
correct prediction of strategies does not determine relative node
probabilities. Additional counterfactual beliefs are required.

At Matching Pennies' equilibrium a player is indifferent between
both pure moves. Rational-choice maximization alone permits all
mixtures and both pure moves against the fixed fair opponent law;
it does not cause a fair randomizer to operate. Interpreting the
equilibrium weights as subjective forecasts can be internally
consistent without matching an empirical frequency prediction.
The same issue occurs when a nonstrict pure equilibrium predicts one
specific move from several tied optima. This is an interpretive
objection to a prediction claim, not a failure of the Nash inequalities
\parencite[pp.~538--539]{AlbertKliemt2021ClassicalGameTheory}.

\section{Normal, extensive and repeated objects remain separate}

The existing finite strategic chapter supplies normal-form existence,
best-response geometry, minimax, correlation and finite approximation.
The present source supplies trees, information sets, off-path plans,
memory and explicit commitment. Normal form preserves payoffs of whole
plans, but does not by itself state which histories form proper subgames.
The mixed--behavioral bridge is conditional on perfect recall and
preserves terminal laws, not a canonical belief at every null history.

A repeated game is another extensive construction: the same stage
game recurs and a plan can condition on observed earlier play. It is
not merely another mixed strategy in the original one-shot matrix.
Albert and Kliemt's selected chapter develops no repeated-game theorem.
The following limited comparison is editorial orientation, using its
actual Prisoner's Dilemma payoffs and explicit added assumptions.

\begin{hbtheorem}{thm:p17-rat91-repetition}{proved=annex:p17.rat91.repetition}{Finite versus infinite repetition under stated observation and payoff assumptions}
Repeat the source Prisoner's Dilemma \(T\geq1\) times with perfect
observation of previous stage profiles, simultaneous moves within
each stage, and total payoff the sum of stage payoffs. Its unique
subgame-perfect behavior is \((D,R)\) at every history of every stage.
At \(T=2\), unilateral first-stage cooperation loses 1 immediately,
and cannot change the equilibrium continuation payoff 1.

Instead allow infinitely many stages with normalized discounted
payoff \((1-\delta)\sum_{t=0}^{\infty}\delta^t u_t\),
\(0\leq\delta<1\). The grim-trigger profile cooperates \((U,L)\)
until any defection is observed, then plays \((D,R)\) forever.
It is subgame perfect exactly when \(\delta\geq1/3\).
At a cooperative history staying gives 3; the best first defection
followed by punishment gives \(4-3\delta\). This result depends
on the infinite horizon, the discount rule, public perfect monitoring
and that specified punishment. It is not a general folk theorem.
\end{hbtheorem}

\section{Reasoning about knowledge and responsible use}

The source's last section starts with a conditional normative point:
if everyone accepts a theory's recommended profile and no unilateral
alternative does better while the other recommendations stand, that
profile is Nash. Such stability does not establish that a recommendation
is uniquely selected, morally justified or accurately forecast.
The chapter describes a tradition seeking correct beliefs from
unlimited rational reasoning and knowledge of the game and theory,
then contrasts its own approach with common-knowledge formulations.
It argues that logical descriptions of ideal knowledge alone are
unlikely to explain real interaction, attributing a sharp criticism
to Selten, while allowing applications supported by independent
psychological hypotheses. These are the authors' methodological
arguments, not experimental estimates of human performance
\parencite[pp.~540--541]{AlbertKliemt2021ClassicalGameTheory}.

The comparison with \cref{sec:p17-epistemic-game-theory} is precise:
that finite type model allows erroneous conjectures and correlated
opponent beliefs; common belief in rationality yields its elimination
result but need not select Nash. The source's correct-prediction
interpretation is a different premise. Likewise,
\cref{thm:p17-ab-two-player,thm:p17-ab-many-player} state actual
conditions on payoff knowledge, rationality, full conjectures and,
in the latter, a common prior with a positive-probability event.
Those theorems distinguish conjectures from actual actions and
null states. The present criticism does not refute them, and those
theorems do not establish their premises for observed people.
The forward rule in \cref{thm:p17-rat91-outside} also differs from
the explicit belief-revision rule of
\cref{thm:p17-egt-reny-reasoning}; agreement of some outcomes is
not identity of off-path reasoning.

Evolutionary models and behavioral theories can use the same game
language with different learning mechanisms, motivations or choice
rules. The chapter points to them as alternatives, without proving
that every such learning process converges to a classical equilibrium.
Its closing acknowledgment credits comments on the exposition;
its references provide routes to historical, philosophical and
mathematical treatments rather than enlarging this presentation's
allocated scope automatically.

\begin{hbmethod}{def:p17-rat91-use}{Choosing and checking the right strategic object}
Specify the terminal histories, player roles, information-set partition,
move labels, payoff interpretation and Nature probabilities. Check
recall before converting mixed to behavioral strategies. Enumerate
complete pure plans to obtain the normal form and apply best-response
or support tests, retaining ties. Then identify complete subgames,
solve them recursively and verify deviations both on and off path.
At an unreached nonsingleton information set, report missing node
beliefs rather than conditioning on zero.

For an application, state why preferences, beliefs, observation,
commitment and time structure describe the proposed situation. Return
the equilibrium profile, its terminal law and its epistemic or
psychological assumptions separately. A successful computation need
not be a successful forecast. Failure cases include an unmodeled
commitment, forgotten own moves, unsupported independent randomization,
or an infinite repetition conclusion transferred to a known finite horizon.
\end{hbmethod}
